\section{Introduction}
\label{sec:Introduction}

In early 2023, engineers at Samsung inadvertently exposed proprietary source code and internal documents by entering them into publicly accessible generative artificial intelligence (AI) tools, prompting the company to ban the use of ChatGPT and similar systems among employees to prevent further leaks \autocite{ray2023samsung}. Security firms now report hundreds of generative-AI-related data violations each month across large organizations as employees routinely engage with AI tools outside approved workflows, risking unintentional disclosure of sensitive information and intellectual property \autocite{arora2025understanding}. At the same time, authoritative cybersecurity assessments warn that advanced AI models can automate and accelerate cyberattacks, lowering the barrier to entry for sophisticated threats and enabling malicious actors to scale operations in ways previously constrained by human and technical resources \autocite{arif2024overview}.\\

These developments illustrate a new class of organizational risk: AI-enabled exposures and attack vectors that propagate faster than traditional defenses can detect or govern, and that bypass established security and compliance controls altogether. AI is not alone. Similar dynamics arise with quantum computing, which threatens the durability of cryptographic infrastructures \autocite{lindsay2020demystifying}; internet of things (IoT) and operational technology convergence, which enables cascading cyber-physical failures \autocite{fu2025systematic}; and software supply-chain platforms, where low-cost, scalable exploits can rapidly propagate across ecosystems \autocite{shen2024understanding}.\\

Against this backdrop, resilience has become the digital regulator’s priority \autocite{dupont2019cyberresilience}. As crises such as COVID-19, global cyberattacks, and systemic information technology (IT) failures have exposed the vulnerabilities of digital infrastructures, resilience has shifted from an abstract aspiration to a core regulatory mandate \autocite{almeida2020challenges}. From financial services to healthcare to critical infrastructure, organizations are now expected to anticipate, absorb, and adapt to disruption -- not only to ensure operational continuity but also to maintain legitimacy in a hyperconnected world \autocite{dupont2019cyberresilience, ivanov2022viable}. The growing uncertainty of the cyber landscape, further amplified by rapid AI adoption and the digital acceleration triggered by the pandemic, has fundamentally redefined expectations for digital systems and those who govern them \autocite{he2022building}.\\

A pronounced mismatch persists between regulatory ambition and organizational capability: most organizations struggle to define what resilience means in a digital context, let alone how to operationalize or measure it \autocite{bjorck2015cyber}. While the reality for organizations is increasingly of continuous, novel and overlapping threats and disruptions, organizational structures constraint them in a way that they can only afford to prepare thoroughly scenarii for one or a few exceptional shocks. One root cause lies in static protocols, rigid control frameworks, and audit-driven compliance \autocite{power1997audit, woods2015four}. These practices -- designed for episodic crises -- provide the appearance of control but do little to sustain adaptive performance in environments of persistent uncertainty, an illusion reinforced by audit culture and decision-making under uncertainty \autocite{power1997audit, satinover2007illusion}. Yet, there is a deep sense of dissonance in the mind of practitioners that their organization may not be capable of preparing fast enough for a greatly accelerated future of technological development and dangers that come along \autocite{power1997audit, hollnagel2013resilience}.\\

This discomfort lies at the heart of the resilience implementation gap. Although resilience is widely invoked in standards, risk frameworks, and regulatory guidelines, it is rarely enacted as a lived organizational capability. Organizations continue to default to compliance-based approaches that signal resilience to external audiences while leaving internal practices unchanged \autocite{holling1996engineering, power1997audit, walker2004resilience}. Empirical evidence confirms this pattern. To get a better view of this paradox in practice, we conducted a ligthweight investigation with 18 cybersecurity and risk professionals. 
%While regulatory expansion through instruments such as DORA, NIS2, the Cyber Resilience Act, or FINMA’s Circular 2023/1 has become the primary driver of initiatives, these efforts often translate into audit preparation and budget compliance rather than into adaptive routines, cross-departmental learning, or proactive culture change.\\  
%To illustrate these tensions, Table \ref{tab:RotaResult} synthesizes recurrent themes emerging from both the literature and our field evidence. 
Table \ref{tab:RotaResult} highlights six patterns: (i) the persistence of ambiguities in the definition of resilience and further in how resilience could be quantified, (ii) the dominance of regulatory pressure and symbolic compliance \autocite{power1997audit}, (iii) the endurance of siloed organizational structures, (iv) the underutilization of learning from incidents, (v) the uneven role of human and cultural factors, and (vi) the paradox between the regulatory injunction for flexibility and organizational demand for precision (see Appendix \ref{appendix:C} for further results and methodological details). Together, these patterns capture the practical contradictions that prevent resilience from being enacted as an organizational capability. Rather than treating resilience as a singular construct, they reveal it as a fragmented phenomenon, somewhat contested within the organization, and surely invoked widely but inconsistently translated into practice.\\


\begin{table}[H]
    \centering
    \singlespacing
    \begin{tabular}{p{0.25\textwidth} p{0.35\textwidth} p{0.35\textwidth}}
        \toprule
        \makecell[tl]{\textbf{Theme from Field} \\  \textbf{Evidence}} & \textbf{Dimension of Resilience in Organization}  &  \textbf{Implications for Practice}\\
        \midrule
        \makecell[tl]{\textbf{Quantification \&} \\ \textbf{Ambiguity}} & \makecell[tl]{Organizations lack operational\\ metrics; resilience remains rhe- \\torical and fragmented across \\ units.} & \makecell[l]{Shared definitions and standar- \\ dized constructs are needed to \\ enable measurement.}\\
         \makecell[tl]{\textbf{Regulatory Pressure} \\ \textbf{\&} \\ \textbf{Symbolic Compliance}} &
        \makecell[tl]{Regulatory audits dominate, \\ producing “tick-box” resilience \\ without adaptive capacity.} &
        \makecell[tl]{Tools must assess adaptive be- \\havior, not just documentation.} \\
        \makecell[tl]{\textbf{Coordination \&} \\ \textbf{Silos}} & \makecell[tl]{IT, risk, compliance units oper- \\ate in silos; limited cross-func- \\tional coordination.} & \makecell[l]{Organizational design should \\ support cross-unit collaboration \\ under crisis.}\\
        \makecell[tl]{\textbf{Learning from} \\ \textbf{Incidents}} & \makecell[tl]{Lessons learned rarely institu-\\tionalized; ad-hoc reviews lack \\ persistence.} & \makecell[l]{Mechanisms are needed to embed \\ feedback and retain knowledge.}\\[3em]
        \makecell[tl]{\textbf{Human \&} \\ \textbf{Cultural Factors}} & \makecell[tl]{Expertise, simulation exercises,\\ and culture matter, but interpre- \\tations vary across units.} & \makecell[l]{Behavioral and cultural aspects \\should be integrated into resilience\\ frameworks.}\\
        \makecell[tl]{\textbf{Flexibility–Precision} \\ \textbf{Paradox}} & \makecell[tl]{Regulators provide flexible defi-\\nitions; firms request precise \\ ones.} & \makecell[l]{Balance between contextual adap-\\tation and measurable guidance \\ must be achieved.}\\
         \bottomrule
    \end{tabular}
    \caption{Field evidence on how organizations interpret and enact resilience in practice (Method : semi-structured interviews, c.f. Appendix \ref{appendix:C} for further results and methodological aspects).}
    \label{tab:RotaResult}
\end{table}

The challenge is compounded by conceptual ambiguity. As resilience has diffused across disciplines and industries, its meaning has fragmented -- ranging from system up-time to business continuity, regulatory posture, or security architecture. This broadening definition further undermines coherence and accountability \autocite{duchek2020organizational, woods2015four}. Even when frameworks exist, incidents and near-misses are seldom leveraged for organizational learning, leaving reflexivity, decentralization, and coordination underdeveloped \autocite{barasa2018what, francis2014metric}. Standards and guidelines continue to proliferate \autocite{linkov2019fundamental}, but organization-ready methods to quantify, simulate, and design resilience remain missing \autocite{kuypers2018designing}. In practice, this means resilience risks remaining an attractive principle with little actual traction on the ground.\\
\\

These insights reinforce the central motivation of this work: resilience is not absent from organizational discourse but enacted in ways that are partial, instrumental, and misaligned with its systemic intent. Efforts remain largely driven by external pressures rather than by endogenous capacity building, leaving organizations without a clear, actionable definition of what it means to build resilience.\\ 
\\

These insights reinforce the central motivation of this work: resilience is not absent from organizational discourse but enacted in ways that are partial, instrumental, and misaligned with its systemic intent. Efforts remain largely driven by external pressures rather than by endogenous capacity building, leaving organizations without a clear, actionable understanding of what it means to build resilience. At its core, resilience requires organizations to operate under uncertainty, adapt continuously, and accept limits to centralized control -- conditions that are organizationally costly, culturally resisted, and therefore difficult to institutionalize. A more robust perspective therefore situates resilience as a socio-technical capability: distributed across actors, sustained through coordination and trust, and enabled by adaptive governance arrangements rather than static control mechanisms.\\
\\

To address this impasse, the paper adopts a deliberately integrative and problem-driven approach. Rather than offering yet another definition of resilience or prescribing best practices, we examine why resilience repeatedly collapses into compliance, symbolism, or fragmentation once translated into organizational settings. We treat ambiguity and paradox not as conceptual failures to be eliminated upfront, but as empirical signals of deeper socio-technical tensions between regulation and practice, control and adaptation, and measurement and learning. By systematically tracing how resilience has been conceptualized, operationalized, and enacted across disciplines -- and where these translations break down in practice -- we shift the focus from what resilience should mean to how it can be rendered observable, comparable, and cumulative without undermining its adaptive nature.\\
\\

Our contribution is to establish resilience as a measurable, socio-technical capability central to the IS field. We move beyond dominant treatments of resilience as rhetoric or regulatory compliance by synthesizing cross-disciplinary insights and organizational evidence to show why resilience consistently fails to materialize in practice. We identify the core conceptual, methodological, and empirical gaps that impede operationalization, propose an integrative framework grounded in behavioral indicators, simulation, and institutional dynamics, and advance a research agenda that transforms resilience from an aspirational label into a tractable theoretical and empirical construct. In doing so, this paper positions resilience as a key frontier for IS scholarship on organizational adaptation under persistent digital disruption.\\
\\

The paper proceeds as follows. Section \ref{sec:Background} synthesizes how resilience has been conceptualized across disciplines and within management and information system research, highlighting how its diffusion has generated persistent conceptual ambiguity. Section \ref{sec:RG} identifies the resulting conceptual, methodological, and empirical gaps that explain why resilience is widely proclaimed yet rarely enacted in practice. Section \ref{sec:RD} advances a research agenda aimed at overcoming this impasse by developing measurable indicators, simulation-based methods, and integrative frameworks that render resilience observable and tractable. Section \ref{sec:ImplicationsResearch} discusses the implications of this agenda for both IS scholarship and organizational practice, emphasizing how resilience can be cultivated as a genuine capability rather than displayed as a compliance artifact. Section \ref{sec:Conclusion} concludes by restating the contribution of the review and outlining pathways toward cumulative and actionable resilience research in the IS field.\\



